% GENERATED FILE. Edit canonical lessons, then run npm run generate:books.
% canonical-source-hash: fnv1a64:833975fba0664c37
% canonical-lessons: AR-C130-tin, AR-C130-bittikh, AR-C130-mishmish, AR-C130-hasa, AR-C130-kaka

\chapter{Figs, Watermelon, Apricots, Soup, Cake}
\label{ch:ar-figs-watermelon-apricots-soup-cake}
\hlchaptermodality{130}

\begin{chapteropening}
\textbf{By the end of this chapter:} \emph{I can say figs, a watermelon, apricots, soup and a cake.}

Complete the last lesson of chapter 130: I can say figs, a watermelon, apricots, soup and a cake.
\end{chapteropening}

\section[\texorpdfstring{\ar{تين} (tīn)}{تين (tīn)}]{\ar{تين} (tīn) — figs}

\label{lesson:AR-C130-tin}



\begin{quote}
Before the new one: say the Arabic for apples, then the Arabic for dates.
\end{quote}

\subsection*{You'll want to know: \ar{تين}}
\textbf{\ar{تين}} — \emph{tīn} — \textquotedblleft{}figs\textquotedblright{}.

\begin{grammarlens}[title={What you've built}]
One more word for this chapter.
\end{grammarlens}

\subsection*{Your turn}
\begin{itemize}
\raggedright
  \item \emph{Say it:} \emph{tīn}
  \item \emph{Say it:} \emph{tīn}, once more
  \item \emph{Say it:} say \emph{tīn}
  \item \emph{Recall:} say the Arabic for apples, then the Arabic for dates, then say \emph{tīn} again
\end{itemize}

\begin{tcolorbox}[breakable,colback=teal!4,colframe=teal!35!black,title={Before you move on}]
What does \ar{تين} mean? (\textquotedblleft{}Figs\textquotedblright{}.) Say it once more.
\end{tcolorbox}
\section[\texorpdfstring{\ar{بطيخ} (biṭṭīkh)}{بطيخ (biṭṭīkh)}]{\ar{بطيخ} (biṭṭīkh) — a watermelon}

\label{lesson:AR-C130-bittikh}



\begin{quote}
Before the new one: say the Arabic for dates, then the Arabic for figs.
\end{quote}

\subsection*{You'll want to know: \ar{بطيخ}}
\textbf{\ar{بطيخ}} — \emph{biṭṭīkh} — \textquotedblleft{}a watermelon\textquotedblright{}.

\begin{grammarlens}[title={What you've built}]
One more word for this chapter.
\end{grammarlens}

\subsection*{Your turn}
\begin{itemize}
\raggedright
  \item \emph{Say it:} \emph{biṭṭīkh}
  \item \emph{Say it:} \emph{biṭṭīkh}, once more
  \item \emph{Say it:} say \emph{biṭṭīkh}
  \item \emph{Recall:} say the Arabic for dates, then the Arabic for figs, then say \emph{biṭṭīkh} again
\end{itemize}

\begin{tcolorbox}[breakable,colback=teal!4,colframe=teal!35!black,title={Before you move on}]
What does \ar{بطيخ} mean? (\textquotedblleft{}A watermelon\textquotedblright{}.) Say it once more.
\end{tcolorbox}
\section[\texorpdfstring{\ar{مشمش} (mishmish)}{مشمش (mishmish)}]{\ar{مشمش} (mishmish) — apricots}

\label{lesson:AR-C130-mishmish}



\begin{quote}
Before the new one: say the Arabic for figs, then the Arabic for a watermelon.
\end{quote}

\subsection*{You'll want to know: \ar{مشمش}}
\textbf{\ar{مشمش}} — \emph{mishmish} — \textquotedblleft{}apricots\textquotedblright{}.

\begin{grammarlens}[title={What you've built}]
One more word for this chapter.
\end{grammarlens}

\subsection*{Your turn}
\begin{itemize}
\raggedright
  \item \emph{Say it:} \emph{mishmish}
  \item \emph{Say it:} \emph{mishmish}, once more
  \item \emph{Say it:} say \emph{mishmish}
  \item \emph{Recall:} say the Arabic for figs, then the Arabic for a watermelon, then say \emph{mishmish} again
\end{itemize}

\begin{tcolorbox}[breakable,colback=teal!4,colframe=teal!35!black,title={Before you move on}]
What does \ar{مشمش} mean? (\textquotedblleft{}Apricots\textquotedblright{}.) Say it once more.
\end{tcolorbox}
\section[\texorpdfstring{\ar{حساء} (ḥasāʾ)}{حساء (ḥasāʾ)}]{\ar{حساء} (ḥasāʾ) — soup}

\label{lesson:AR-C130-hasa}



\begin{quote}
Before the new one: say the Arabic for a watermelon, then the Arabic for apricots.
\end{quote}

\subsection*{You'll want to know: \ar{حساء}}
\textbf{\ar{حساء}} — \emph{ḥasāʾ} — \textquotedblleft{}soup\textquotedblright{}.

\begin{grammarlens}[title={What you've built}]
One more word for this chapter.
\end{grammarlens}

\subsection*{Your turn}
\begin{itemize}
\raggedright
  \item \emph{Say it:} \emph{ḥasāʾ}
  \item \emph{Say it:} \emph{ḥasāʾ}, once more
  \item \emph{Say it:} say \emph{ḥasāʾ}
  \item \emph{Recall:} say the Arabic for a watermelon, then the Arabic for apricots, then say \emph{ḥasāʾ} again
\end{itemize}

\begin{tcolorbox}[breakable,colback=teal!4,colframe=teal!35!black,title={Before you move on}]
What does \ar{حساء} mean? (\textquotedblleft{}Soup\textquotedblright{}.) Say it once more.
\end{tcolorbox}
\section[\texorpdfstring{\ar{كعكة} (kaʿka)}{كعكة (kaʿka)}]{\ar{كعكة} (kaʿka) — a cake}

\label{lesson:AR-C130-kaka}



\begin{quote}
Before the new one: say the Arabic for apricots, then the Arabic for soup.
\end{quote}

\subsection*{You'll want to know: \ar{كعكة}}
\textbf{\ar{كعكة}} — \emph{kaʿka} — \textquotedblleft{}a cake\textquotedblright{}.

\begin{grammarlens}[title={What you've built}]
One more word for this chapter.
\end{grammarlens}

\subsection*{Your turn}
\begin{itemize}
\raggedright
  \item \emph{Say it:} \emph{kaʿka}
  \item \emph{Say it:} \emph{kaʿka}, once more
  \item \emph{Say it:} say \emph{kaʿka}
  \item \emph{Recall:} say the Arabic for apricots, then the Arabic for soup, then say \emph{kaʿka} again
\end{itemize}

\begin{tcolorbox}[breakable,colback=teal!4,colframe=teal!35!black,title={Before you move on}]
What does \ar{كعكة} mean? (\textquotedblleft{}A cake\textquotedblright{}.) Say it once more.
\end{tcolorbox}
